\section{How the cancellations were identified}
Division by $x^2$ near $x=0$ revealed a difficulty that remained hidden in the interior. Two
independent evaluations of a formula could agree if both omitted the same incidence. The discrepancy
with the expected slope led us to vary the square weights and count the physical lines one by one.
When two squares share another line, its vacancy enters only once. Applying this rule in both
orders of the product gave the corrections $\Delta_1,\Delta_2$. The catalog in the supplement
documents the blind comparison of implementations, the line-by-line evaluator, and the use of
profile deformations to locate the two omissions.

\paragraph{How to retain the sign when summing the errors.}
A bound on $|R_Q(n)|$ allows $nR_Q(n)$ to oscillate. To retain the sign we compare, as in
Section~\ref{sec:signedmap}, the probabilities of the same stopped history under the two laws.
At arrival there are groups whose size grows with $n$. Conditioning on the rows involved allows
us to bound the sum of these groups in mean, with a common constant for all small cut-offs. This
is the estimate the proof needs. Pointwise control on each board is a stronger requirement.

\paragraph{From the permanent of the mask to the error of a pair.}
Occupying the pair removes two rows and two columns. Thus it matters how much the error of the
permanent changes on passing to the minor. Divoux and Wang give the leading determinant in the
diffuse class, McCullagh develops formal corrections~\cite{dw,mcc}, and Li treats perturbations
of sublinear line norm~\cite{li}. Estimate $(L^*)$ selects the first excess class of the Gaussian
expansion and controls the remainder uniformly. To obtain $J$ the availability of the other two
vertices is held fixed while the weights of the pair are varied. Rank-two reinsertion identifies
the response of a single line even when both diagonals are available. Thus the square interaction
contributes only once to $-S$.

\paragraph{From an algebraic check to a consequence of the theorem.}
The cancellation of the linear term at the endpoint was checked by computer in a system of 42
transcribed pieces \status{O}. The proof for the physical profile follows a different argument:
two cut-offs approximate the same initial quantity and give the Cauchy criterion. A nonzero slope
$k'(0)$ would produce a logarithmic divergence. Convergence of the truncated integral forces
$k'(0)=0$ (Appendix~\ref{app:closing}).

\paragraph{From the formula to a number.}
The first quadratures, based on a few central squares, gave errors of $10^{-5}$ to $10^{-3}$ in
individual contributions. The rectangle method, adapted from~\cite[\S4]{tc}, sums all the centers
of each pattern on the same grid. This rearrangement is exact. The passage from the grid to the
continuum is estimated by extrapolation \status{N}. Refining the profile made it possible to enclose
it with radii below $2.1\cdot10^{-15}$ and obtain an interval for $Q_2$ at one real node. The trace
$S_{13}$ explains why the interval is still wide: bounding the entries of the row and column block
separately loses their dependence. Section~\ref{sec:numcert} quantifies this limitation.
